
This study provides evidence that RLHF instruction tuning degrades the discriminative quality of confidence signals in large language models. Across two model families (Qwen and Mistral), instruction-tuned models show reduced AUROC for margin-based correctness prediction (mean decrease: 3.03 percentage points), disproportionate margin inflation for incorrect predictions (3-17x), percentile-normalized slope attenuation (34-40\%), and Brier refinement degradation (0.02-0.05).

These findings characterize the distortion as geometric rather than scalar: RLHF reshapes the confidence-correctness relationship rather than merely rescaling probabilities. This distinction suggests that post-hoc calibration methods designed for scalar miscalibration may not fully address discriminative degradation in RLHF-tuned models.

Future work should: (1) validate findings across additional model families including Llama, (2) test across model scales beyond 7B parameters, (3) empirically verify whether temperature scaling improves ECE while leaving AUROC degraded, (4) investigate domain-specific calibration patterns, and (5) explore calibration-aware RLHF training procedures that preserve discriminative quality.
